\chapter{Discussion}

\section{The DRAM onset below nominal L3}

The latency staircase steps up to DRAM near an 8 to 16 MB working set, while the
L3 is 33 MB. Single threaded random pointer chasing does not reach the full L3
before latency climbs, because the random 64 byte stride thrashes the
translation lookaside buffer, whose reach with 4 KB pages is limited, and it
stresses L3 associativity and replacement. Effective capacity under this access
pattern is roughly half the nominal size. This is expected
behavior~\cite{drepper2007}, and the sysfs cross check is designed to flag it
rather than smooth it away. Large pages would isolate a purer DRAM latency, at
the cost of measuring something less representative of ordinary code.

\section{The bandwidth wall}

Bandwidth saturates at four to five threads. This is the defining feature of a
memory bound workload on a dual channel desktop: the DRAM channels are the
bottleneck, and once they are full, more cores cannot help. The practical
consequence is that the P versus E core question is moot for STREAM, since the
knee arrives long before the core count does. It would matter for a compute
bound workload, where the E cores add throughput, which is visible in the all
core FLOPS number.

\section{Why the tile prediction misses}

The BLIS analytical model assumes a packed micro kernel in which the L1 resident
panel of $B$ is only $N_r$ elements wide. The validation kernel here is
deliberately naive: it does not pack, and its inner loop spans the full width of
$B$. Each reused row of $B$ is therefore the full matrix width, several
kilobytes, so only a handful of rows fit in L1 or L2 at once, and the optimum
moves to a much smaller $K_c$ than the packed model wants. The prediction is not
wrong for the kernel it describes; it is describing a different kernel. Closing
the gap means implementing the packed micro kernel the model assumes, which is
the natural next step and is recorded as future work.

\section{Limitations}

The clock is measured in software, not read from a counter, so it carries a few
percent of uncertainty; this is why efficiencies are reported as approximate and
a match to within one percent is treated as confirmation rather than a precise
claim. Bandwidth uses float arrays and the classical STREAM byte convention. The
hybrid P and E identification is heuristic under WSL2. None of these affect the
qualitative conclusions, and each is stated where it applies.
